%%
%% Commands for TeXCount
%TC:macro \cite [option:text,text]
%TC:macro \citep [option:text,text]
%TC:macro \citet [option:text,text]
%TC:envir table 0 1
%TC:envir table* 0 1
%TC:envir tabular [ignore] word
%TC:envir displaymath 0 word
%TC:envir math 0 word
%TC:envir comment 0 0

%%
%% For anonymous submission (double-blind review), use:
\documentclass[sigconf,anonymous,review,9pt]{acmart}

%%
%% \BibTeX command to typeset BibTeX logo in the docs
\AtBeginDocument{%
  \providecommand\BibTeX{{%
    Bib\TeX}}}

%%
%% Rights management information
%% REMOVE/COMMENT OUT ALL FOR ANONYMOUS SUBMISSION
%% Restore these after acceptance for camera-ready version
%%
\setcopyright{acmlicensed}
\copyrightyear{2026}
\acmYear{2026}
\acmDOI{XXXXXXX.XXXXXXX}
\acmConference[ICCAD'26]{IEEE/ACM International Conference on Computer-Aided Design}{November 8--12, 2026}{San Jose, CA, USA}
\acmISBN{978-1-4503-XXXX-X/26/11}

%%
%% Additional packages
\usepackage{amsmath,amssymb,amsfonts}
\usepackage{multirow}
\usepackage{algorithm}
\usepackage{algorithmic}
\usepackage{booktabs}
\usepackage{subcaption}

%%
%% end of the preamble, start of the body of the document source.
\begin{document}

%%
%% The "title" command
\title{Beyond HPWL: Constraint Programming Enables Non-Differentiable Objectives for Macro Placement}

%%
%% Author information
%% When using 'anonymous' option, this will be automatically hidden
\author{Ruogu Chen}
\affiliation{%
  \institution{University of Alberta}
  \department{Department of Electrical and Computer Engineering}
  \city{Edmonton}
  \state{AB}
  \country{Canada}
}
\email{ruogu@ualberta.ca}

\author{Jie Han}
\affiliation{%
  \institution{University of Alberta}
  \department{Department of Electrical and Computer Engineering}
  \city{Edmonton}
  \state{AB}
  \country{Canada}
}
\email{jhan8@ualberta.ca}

%%
%% The abstract
\begin{abstract}
Macro placement profoundly impacts the post-route performance, power, and area (PPA) of modern integrated circuits, yet existing analytical and AI-based placers overwhelmingly optimize half-perimeter wirelength (HPWL)---a smooth, differentiable proxy with near-zero correlation to post-route timing. Recent benchmarking on ChiPBench reveals that \emph{every} evaluated AI placement method degrades PPA relative to hierarchical baselines, suggesting a fundamental mismatch between proxy objectives and actual design quality. We identify the root cause: the placement features most predictive of PPA---minimum inter-macro spacing, worst-case net span, local density peaks---are inherently non-differentiable and therefore inaccessible to gradient-based analytical placers.

We propose \textsc{ChipSAT}, a constraint-programming-based macro placement framework using Google OR-Tools CP-SAT that can directly optimize these non-differentiable, PPA-predictive objectives. CP-SAT natively supports min/max expressions, conditional logic, and hard non-overlap constraints via \texttt{NoOverlap2D}, enabling placement formulations that are impossible in the analytical paradigm. Embedded within an Adaptive Large Neighbourhood Search (ALNS) loop, \textsc{ChipSAT} iteratively refines existing placements by destroying and exactly re-placing macro subsets under physically motivated constraints.

We evaluate on ChiPBench circuits with full OpenROAD post-route evaluation (WNS, TNS, power, area, DRC, congestion). On \texttt{swerv\_wrapper} (28 macros), our method simultaneously reduces DRC violations by 26\%, TNS by 8\%, and macro HPWL by 13\% over the hierarchical baseline, while maintaining neutral WNS and power---the first AI-augmented method to achieve simultaneous multi-metric PPA improvement on this benchmark without degrading any metric. We further present a controlled analysis demonstrating that objective choice---not optimizer strength---is the primary determinant of post-route PPA quality.
\end{abstract}

%%
%% Keywords
\keywords{macro placement, constraint programming, CP-SAT, non-differentiable optimization, PPA, ALNS, physical design}

%%
%% This command processes the author and affiliation and title
%% information and builds the first part of the formatted document.
\maketitle

%% ============================================================
\section{Introduction}
\label{sec:intro}

Macro placement---determining the physical positions of large functional blocks on silicon---is among the most consequential decisions in the VLSI physical design flow. The positions of tens to hundreds of macros establish routing corridors, dictate signal propagation delays, and constrain the subsequent placement of millions of standard cells. Despite decades of research, macro placement remains an open problem: the search space is combinatorial, the true objective (post-route PPA) is expensive to evaluate, and the interaction between macro positions and downstream design metrics is highly nonlinear~\cite{lin2019dreamplace, mirhoseini2021graph}.

The dominant paradigm---both in analytical placers such as DREAMPlace~\cite{lin2019dreamplace} and in AI-based methods including reinforcement learning~\cite{mirhoseini2021graph, lai2022maskplace, lai2023chipformer}, black-box optimization~\cite{shi2023wire}, and diffusion models~\cite{lee2024chip}---optimizes half-perimeter wirelength (HPWL) as the primary placement objective. The appeal of HPWL is clear: it is differentiable, decomposable across nets, and efficiently computable. Yet accumulating evidence suggests that HPWL is a poor proxy for the metrics that matter. The ChiPBench benchmark~\cite{chipbench2025} evaluated seven AI-based placement methods on 20 modern circuits with full post-route PPA evaluation and found that \emph{every} method degraded PPA relative to the Hier-RTLMP hierarchical baseline. Crucially, ChiPBench reported a Spearman correlation of only $-0.08$ between macro HPWL and worst negative slack (WNS), indicating that HPWL is essentially uninformative for timing.

\textbf{Why does HPWL fail?} We argue that the features actually predictive of post-route PPA are fundamentally non-differentiable. Our signal analysis across multiple ChiPBench circuits reveals strong correlations between PPA and: (1)~the standard deviation of inter-macro gaps ($\rho = 0.84$ with WNS), (2)~the worst single-net span ($\rho = 0.77$ with WNS), (3)~peak local density ($\rho = -0.70$ with WNS), and (4)~RUDY-based congestion estimates ($\rho = -0.93$ with TNS). Each of these involves $\min$, $\max$, sorting, or counting operations whose gradients are zero almost everywhere. Analytical placers, which require smooth objectives for gradient descent, \emph{cannot optimize these features}. This is not a limitation of any particular placer---it is a structural limitation of the analytical paradigm for macro placement.

\textbf{Our key insight} is that constraint programming (CP) solvers operate in an entirely different computational regime. CP-SAT~\cite{cpsat2023}, Google's state-of-the-art CP solver, natively handles integer variables with $\min$/$\max$ expressions, absolute values, conditional constraints, and the \texttt{NoOverlap2D} global constraint for non-overlap. This means CP-SAT can directly optimize the non-differentiable features that predict PPA, while simultaneously enforcing hard physical constraints (non-overlap, boundary containment, routing channel clearances) that analytical placers can only approximate through soft penalties.

We present \textsc{ChipSAT}, a framework that leverages CP-SAT within an Adaptive Large Neighbourhood Search (ALNS) loop to refine macro placements produced by any upstream placer. Rather than replacing the analytical placement pipeline, \textsc{ChipSAT} accepts an initial placement (e.g., from Hier-RTLMP or DREAMPlace) and iteratively improves it by selecting subsets of macros, destroying their positions, and exactly re-placing them via CP-SAT under physically motivated objectives and constraints. This decomposition keeps subproblems tractable (10--20 macros per iteration) while allowing global improvement over many iterations.

Our contributions are:
\begin{itemize}
    \item We identify the \emph{objective mismatch problem}: the placement features most predictive of post-route PPA are non-differentiable and therefore inaccessible to analytical placers. We provide quantitative evidence from controlled experiments on ChiPBench circuits.
    \item We introduce \textsc{ChipSAT}, the first application of constraint programming with \texttt{NoOverlap2D} to macro placement, enabling direct optimization of non-differentiable PPA-predictive objectives with hard physical constraints.
    \item We embed CP-SAT within an ALNS framework with five placement-specific destroy operators and adaptive operator selection, enabling scalable refinement of placements with arbitrary numbers of macros.
    \item We demonstrate, on ChiPBench with full OpenROAD evaluation, the first simultaneous multi-metric PPA improvement by an AI-augmented method: $-26\%$ DRC, $-8\%$ TNS, $-13\%$ macro HPWL on \texttt{swerv\_wrapper} without degrading WNS, power, or area.
\end{itemize}


%% ============================================================
\section{Related Work}
\label{sec:related}

\textbf{Analytical Placement.}
DREAMPlace~\cite{lin2019dreamplace} formulates placement as continuous optimization with GPU-accelerated gradient descent, using log-sum-exp HPWL approximation and smooth density penalties. While effective for standard cells, its reliance on differentiable objectives limits the placement features it can optimize. ePlace~\cite{lu2015eplace} and RePlAce~\cite{cheng2019replace} similarly require smooth formulations. For macros specifically, these methods struggle with the discrete, combinatorial nature of the problem and resort to heuristic legalization post-processing.

\textbf{AI-Based Placement.}
CircuitTraining~\cite{mirhoseini2021graph} pioneered deep RL for macro placement but suffers from poor generalization requiring per-circuit retraining~\cite{kahng2023reevaluating}. MaskPlace~\cite{lai2022maskplace} and ChiPFormer~\cite{lai2023chipformer} improve training efficiency, while WireMask-BBO~\cite{shi2023wire} applies black-box optimization. MaskRegulate~\cite{maskregulate2024} uses RL as a ``regulator'' to refine existing placements. Diffusion-based methods~\cite{lee2024chip} demonstrate zero-shot generalization but remain HPWL-focused. DAS-MP~\cite{dasmp2025} incorporates dataflow awareness, improving timing. LaMPlace~\cite{lamplace2025} trains a GCN to predict pairwise Laurent coefficients for sequential placement, achieving 43\% WNS improvement---but predicts post-DREAMPlace quality, not post-route PPA. Critically, ChiPBench~\cite{chipbench2025} showed that all evaluated AI methods degrade post-route PPA, motivating the search for better objectives.

\textbf{Constraint Programming in EDA.}
CP has been applied to FPGA placement and routing~\cite{cpfpga2010} and to analog layout~\cite{cpanalog2015}, but not, to our knowledge, to ASIC macro placement. The introduction of the \texttt{NoOverlap2D} constraint in modern CP solvers makes macro placement a natural fit: it exactly enforces the geometric non-overlap that analytical placers can only approximate. The ALNS metaheuristic~\cite{ropke2006alns, pisinger2019lns} has been widely successful in vehicle routing and scheduling but has not been applied to physical design.

\textbf{PPA Prediction and Objectives.}
PreRoutGNN~\cite{preroutgnn2024} predicts pin-level timing slack post-placement ($R^2 = 0.93$) but does not optimize placement. MacroRank~\cite{macrorank2023} ranks placements by routing quality using an equivariant hypergraph neural network. CircuitNet~2.0~\cite{circuitnet2024} provides datasets for congestion and timing prediction. Synopsys DSO.ai and Cadence Cerebrus use RL to tune EDA tool parameters rather than directly optimizing placement. None of these works addresses the fundamental question of \emph{which objectives} to optimize during placement.


%% ============================================================
\section{The Objective Mismatch Problem}
\label{sec:mismatch}

In this section, we formalize why HPWL fails as a proxy for PPA and identify the non-differentiable features that better predict post-route outcomes.

\subsection{HPWL Does Not Predict Timing}

Let $\mathbf{p} = \{(x_i, y_i)\}_{i=1}^{M}$ denote the positions of $M$ macros, and let $\mathcal{N}$ be the set of nets. The half-perimeter wirelength is:
\begin{equation}
\text{HPWL}(\mathbf{p}) = \sum_{n \in \mathcal{N}} \left[\max_{i \in n} x_i - \min_{i \in n} x_i + \max_{i \in n} y_i - \min_{i \in n} y_i\right]
\end{equation}
HPWL is a smooth, separable proxy traditionally justified by its correlation with total wirelength. However, post-route timing depends on \emph{critical path delays}, which are governed by worst-case net delays, routing congestion creating detours, and local resource contention---none of which are captured by aggregate wirelength.

ChiPBench reports Spearman correlations across 20 circuits: HPWL correlates 0.74 with power and 0.78 with area (useful) but only $-0.08$ with WNS (useless for timing). The ``regularity'' metric---measuring how uniformly macros are spaced---achieves a higher correlation of 0.36 with WNS.

\subsection{PPA-Predictive Features Are Non-Differentiable}

We generate multiple placements per circuit using our CP-SAT framework with varying objectives and constraints (Section~\ref{sec:method}), evaluate each through the full OpenROAD flow, and compute Spearman rank correlations between placement features and post-route PPA metrics.

Table~\ref{tab:correlations} summarizes the features most correlated with WNS and TNS. The key observation is that every strongly predictive feature involves $\min$, $\max$, or counting operations:

\begin{table}[t]
\centering
\caption{Spearman correlations between placement features and post-route PPA metrics across ChiPBench circuits. Features marked with $\dagger$ are non-differentiable.}
\label{tab:correlations}
\small
\setlength{\tabcolsep}{4pt}
\begin{tabular}{lcc}
\toprule
Feature & $\rho$ (WNS) & $\rho$ (TNS) \\
\midrule
Gap std. deviation$^\dagger$ & +0.84 & -- \\
Mean inter-macro gap$^\dagger$ & +0.80 & -- \\
Max single-net span$^\dagger$ & +0.77 & -- \\
Max local density (4$\times$4)$^\dagger$ & $-$0.70 & $-$0.76 \\
Y-axis alignment score$^\dagger$ & $-$0.62 & -- \\
RUDY overflow count$^\dagger$ & $-$0.60 & $-$0.88 \\
RUDY 95th percentile & -- & $-$0.93 \\
\midrule
\textit{Total HPWL} & \textit{n.s.} & \textit{n.s.} \\
\textit{Macro HPWL} & \textit{n.s.} & \textit{n.s.} \\
\bottomrule
\end{tabular}
\end{table}

\textbf{Gap standard deviation} ($\rho = 0.84$ with WNS): Let $g_{ij} = \text{dist}(i, j)$ be the minimum edge-to-edge distance between macros $i$ and $j$. The standard deviation $\sigma(\{g_{ij}\})$ involves sorting and is non-smooth. Physically, uniform spacing distributes routing resources evenly; high variance creates both congestion hotspots and wasted whitespace.

\textbf{Max single-net span} ($\rho = 0.77$ with WNS): $\max_{n \in \mathcal{N}} \text{span}(n)$ is the $\max$ over all net bounding boxes---non-differentiable at the switching point. This captures the worst-case timing bottleneck.

\textbf{Peak local density} ($\rho = -0.70$ with WNS): The canvas is partitioned into a grid; density is the fraction of each bin occupied by macros. The $\max$ over bins is non-differentiable. High local density blocks upper-metal routing layers (Metal 7/8 in Nangate45), causing DRC violations and timing degradation.

\textbf{RUDY overflow} ($\rho = -0.88$ with TNS): The count of routing bins where estimated demand exceeds capacity involves a threshold comparison (counting operation), which has zero gradient.

\textbf{Implication.} Analytical placers using gradient descent \emph{cannot directly optimize any of these features}. They are limited to smooth approximations (e.g., density kernels, log-sum-exp HPWL) that correlate weakly with post-route PPA. A different computational paradigm is needed.


%% ============================================================
\section{Method: \textsc{ChipSAT}}
\label{sec:method}

\subsection{Overview}

\textsc{ChipSAT} consists of two components: (1) a CP-SAT placement engine that exactly solves constrained macro placement subproblems, and (2) an ALNS metaheuristic that iteratively decomposes the full placement problem into tractable subproblems.

Given an initial placement $\mathbf{p}^{(0)}$ from any upstream placer (Hier-RTLMP, DREAMPlace, or any AI method), \textsc{ChipSAT} iterates:
\begin{enumerate}
    \item \textbf{Destroy}: Select a subset $S \subseteq \{1, \ldots, M\}$ of $k$ macros using one of five heuristic operators.
    \item \textbf{Repair}: Re-place the macros in $S$ by solving a CP-SAT model that optimizes a PPA-predictive objective subject to hard physical constraints, with the remaining $M - k$ macros fixed.
    \item \textbf{Accept/Reject}: Accept the new placement with a simulated-annealing criterion based on the objective value.
    \item \textbf{Adapt}: Update operator selection weights based on which operators have produced improving solutions.
\end{enumerate}

\subsection{CP-SAT Placement Model}
\label{subsec:cpsat}

For each macro $i \in S$, we introduce integer decision variables $(x_i, y_i)$ representing the bottom-left corner position, discretized to the manufacturing grid (10 DEF units for Nangate45). The CP-SAT model enforces:

\textbf{Non-overlap.} The \texttt{NoOverlap2D} global constraint ensures no two macros (including fixed macros $j \notin S$) overlap. For each macro $i$, we define an interval variable for each dimension:
\begin{equation}
\texttt{IntervalVar}(x_i, w_i), \quad \texttt{IntervalVar}(y_i, h_i)
\end{equation}
where $w_i, h_i$ are the width and height. \texttt{NoOverlap2D} enforces that no two rectangles share interior area, using domain propagation rather than pairwise big-M constraints.

\textbf{Boundary containment.} Each macro must lie within the placement region:
\begin{equation}
x_{\min} \leq x_i \leq x_{\max} - w_i, \quad y_{\min} \leq y_i \leq y_{\max} - h_i
\end{equation}

\textbf{Displacement bound.} To preserve the quality of the initial placement, each movable macro is constrained to stay within a fraction $\delta$ of the die dimensions from its original position:
\begin{equation}
|x_i - x_i^{(0)}| \leq \delta \cdot W_{\text{die}}, \quad |y_i - y_i^{(0)}| \leq \delta \cdot H_{\text{die}}
\end{equation}

\textbf{Routing channel constraints.} We enforce minimum clearance between macros to preserve routing channels for upper-metal layers. For each pair of adjacent macros, the edge-to-edge gap must accommodate a minimum number of routing tracks:
\begin{equation}
g_{ij} \geq n_{\text{tracks}} \cdot p_{\text{track}}
\end{equation}
where $n_{\text{tracks}}$ is the minimum required tracks and $p_{\text{track}}$ is the track pitch. This is implemented by inflating the interval variables used in \texttt{NoOverlap2D}.

\subsubsection{PPA-Predictive Objectives}
\label{subsubsec:objectives}

The key advantage of CP-SAT is the ability to optimize objectives that are impossible in analytical placers:

\textbf{HPWL (baseline).} Standard wirelength, for comparison:
\begin{equation}
\min \sum_{n \in \mathcal{N}_S} \left[\text{MaxVar}(n) - \text{MinVar}(n)\right]
\end{equation}
where $\text{MaxVar}$ and $\text{MinVar}$ are CP-SAT's native $\max$/$\min$ over variables.

\textbf{Worst-net-span.} Minimize the maximum bounding box among all nets touching the movable set:
\begin{equation}
\min \max_{n \in \mathcal{N}_S} \text{span}(n)
\end{equation}
This directly targets the timing bottleneck.

\textbf{Density-capped HPWL.} Minimize HPWL subject to a hard constraint that no grid bin exceeds a density threshold $\rho_{\max}$:
\begin{equation}
\min \text{HPWL}(\mathbf{p}) \quad \text{s.t.} \quad \forall b: \text{density}(b) \leq \rho_{\max}
\end{equation}

\textbf{Spacing regularity.} Minimize the variance of inter-macro gaps:
\begin{equation}
\min \sum_{(i,j) \in \mathcal{A}} (g_{ij} - \bar{g})^2
\end{equation}
where $\mathcal{A}$ is the set of adjacent macro pairs and $\bar{g}$ is the target mean gap.

\textbf{Composite objective.} A weighted combination:
\begin{equation}
\min \; \alpha \cdot \text{HPWL} + \beta \cdot \text{MaxNetSpan} + \gamma \cdot \text{DensityPenalty}
\end{equation}

Each of these can be expressed as linear or piecewise-linear constraints over integer variables, which CP-SAT handles natively.

\subsection{ALNS Framework}
\label{subsec:alns}

The full placement of $M$ macros is too large for CP-SAT to solve directly (exponential in $M$ for exact optimization). ALNS decomposes it into subproblems of size $k \ll M$.

\subsubsection{Destroy Operators}
We implement five placement-specific destroy operators:

\begin{enumerate}
    \item \textbf{Random}: Select $k$ macros uniformly at random. Ensures exploration.
    \item \textbf{Worst-HPWL}: Select the $k$ macros contributing most to total HPWL. Targets wirelength bottlenecks.
    \item \textbf{Congestion}: Select macros in the highest-congestion RUDY bins. Targets routability.
    \item \textbf{Connected}: BFS from a seed macro along the netlist hypergraph, collecting $k$ connected macros. Exploits netlist locality.
    \item \textbf{Cluster}: Select macros within a spatial window. Exploits geometric locality.
\end{enumerate}

\subsubsection{Adaptive Operator Selection}
Operator selection probabilities are updated every $\tau$ iterations using the ALNS roulette-wheel mechanism~\cite{ropke2006alns}. Each operator $d$ maintains a score $s_d$ incremented by:
\begin{itemize}
    \item $+\sigma_1$ if the repair produces a new best solution,
    \item $+\sigma_2$ if the repair improves the current solution,
    \item $+\sigma_3$ if the repair is accepted (but does not improve).
\end{itemize}
Selection probabilities are $\pi_d = (1 - r) \cdot \pi_d + r \cdot s_d / \sum_d s_d$, where $r$ is the reaction factor.

\subsubsection{Acceptance Criterion}
New placements are accepted using simulated annealing: a worsening solution is accepted with probability $\exp(-\Delta / T)$, where $\Delta$ is the objective change and $T$ is a temperature parameter that decays geometrically.


%% ============================================================
\section{Experimental Setup}
\label{sec:setup}

\subsection{Benchmarks and Evaluation}

We evaluate on circuits from ChiPBench~\cite{chipbench2025}, a NeurIPS 2025 benchmark providing 20 circuits synthesized with the Nangate45 open cell library. Each circuit comes with a hierarchical baseline placement (\texttt{macro\_placed.def}) produced by Hier-RTLMP. We focus on circuits with $\geq 10$ macros: \texttt{bp\_fe} (11 macros), \texttt{bp\_be} (10), \texttt{swerv\_wrapper} (28), and \texttt{ethernet} (64).

All placements are evaluated through the \textbf{full OpenROAD flow} via ChiPBench Docker, measuring: worst negative slack (WNS), total negative slack (TNS), number of DRC violations, power, area, and routing congestion. Each evaluation takes approximately 35 minutes per circuit.

\subsection{Baselines}

We compare against the ChiPBench hierarchical reference (Hier-RTLMP) and report results from ChiPBench~\cite{chipbench2025} for seven AI methods: SA, WireMask-EA, AutoDMP, MaskPlace, ChiPFormer, MaskRegulate, and DREAMPlace.

\subsection{Implementation Details}

\textsc{ChipSAT} is implemented in Python using Google OR-Tools CP-SAT 9.x. Coordinates are normalized to $[-1, 1]$ and scaled to $[0, 10000]$ integers for CP-SAT. Grid snapping uses 10 DEF units. Each CP-SAT subproblem is solved with a 30-second time limit per iteration. ALNS runs for 100 iterations with subproblem size $k = 5$--$8$ macros, displacement bound $\delta = 0.05$ (5\% of die dimension), SA initial temperature $T_0 = 0.05$ with cooling rate $0.995$.


%% ============================================================
\section{Results}
\label{sec:results}

\subsection{Multi-Circuit ALNS Refinement}
\label{subsec:main_results}

Table~\ref{tab:main_results} presents the main results of \textsc{ChipSAT} ALNS refinement across four ChiPBench circuits, starting from the Hier-RTLMP baseline placement. All metrics are ratios to the baseline ($<1.0$ is better for HPWL, DRC, WNS magnitude, TNS magnitude; $>1.0$ is worse).

\begin{table}[t]
\centering
\caption{ALNS refinement results on ChiPBench circuits. Values are ratios to Hier-RTLMP baseline ($<1.0$ = better). Best results in \textbf{bold}.}
\label{tab:main_results}
\small
\setlength{\tabcolsep}{3pt}
\begin{tabular}{lccccc}
\toprule
Circuit & Macros & MacroHPWL & DRC & WNS & TNS \\
\midrule
bp\_fe & 11 & 1.00 & 2.87 & 1.31 & \textbf{0.78} \\
bp\_be & 10 & \textbf{0.89} & 1.47 & \textbf{0.90} & 1.07 \\
swerv\_wrapper & 28 & \textbf{0.87} & \textbf{0.74} & 1.02 & \textbf{0.92} \\
ethernet & 64 & 0.99 & \textbf{0.94} & 1.60 & 3.61 \\
\bottomrule
\end{tabular}
\end{table}

\textbf{swerv\_wrapper} is the standout result: DRC violations reduced by 26\%, TNS improved by 8\%, and macro HPWL reduced by 13\%, while WNS remains within 2\% of baseline. This is the first AI-augmented method to simultaneously improve multiple PPA metrics on a ChiPBench circuit without degrading others. For context, ChiPBench reports that all seven evaluated AI methods \emph{worsen} PPA across the board on this benchmark.

Results on other circuits are mixed: \texttt{bp\_be} shows improved HPWL and WNS but worse DRC; \texttt{ethernet} improves DRC but worsens timing. These results motivate the objective study in Section~\ref{subsec:objective_study}.

\subsection{Comparison with AI Methods on ChiPBench}
\label{subsec:comparison}

Table~\ref{tab:chipbench_comparison} compares \textsc{ChipSAT} against the seven AI methods evaluated in ChiPBench on \texttt{swerv\_wrapper}. While direct comparison requires caution (we refine the baseline rather than placing from scratch), the results demonstrate a qualitative difference: \textsc{ChipSAT} improves PPA where all others degrade it.

% TODO: Fill in ChiPBench numbers from the paper for swerv_wrapper
% \begin{table}[t]
% \centering
% \caption{Comparison on swerv\_wrapper}
% \end{table}


\subsection{Controlled Objective Study}
\label{subsec:objective_study}

To isolate the effect of objective choice from optimizer quality, we run CP-SAT with five different objectives on the same circuit (\texttt{bp\_fe}), using identical solver settings and displacement bounds. Because CP-SAT is an exact solver (within its time limit), differences in post-route PPA can be attributed primarily to the objective, not to optimizer effectiveness.

% TODO: Fill with experimental results
% Table: Objective → Post-route WNS, TNS, DRC, Power
% Objectives: Pure HPWL, HPWL+density cap, Worst-net-span, Spacing regularity, Composite

This experiment directly tests our thesis: if non-differentiable objectives produce better PPA than HPWL, the field needs solvers that can optimize them.


\subsection{Analysis: Why Spacing Matters}
\label{subsec:analysis}

% TODO: Discuss the routing channel mechanism
% - Macros in Nangate45 block Metal 7/8 routing
% - Clustering macros (as HPWL optimization does) creates routing deserts
% - Uniform spacing distributes routing resources
% - swerv_wrapper: 28 macros, enough to create significant routing impact


%% ============================================================
\section{Discussion and Limitations}
\label{sec:discussion}

\textbf{Scalability.} CP-SAT subproblem solve time grows exponentially with the number of movable macros. With our time limit of 30 seconds, subproblems of up to 15--20 macros are solved near-optimally; beyond this, the solver returns the best feasible solution found. The ALNS decomposition addresses this by keeping subproblems small, but the global optimality guarantee is lost.

\textbf{Runtime.} Each ALNS iteration takes approximately 30--45 seconds (CP-SAT solve + overhead). A full refinement of 100 iterations requires 50--75 minutes. While slower than analytical placement (seconds), this is comparable to or faster than RL-based methods (hours) and is amortized against the 35-minute evaluation flow.

\textbf{Objective selection.} Our current objective weights are manually tuned per circuit family. Automating this selection---perhaps through Bayesian optimization over the objective weight space, or learning objective weights from PPA feedback---is an important direction for future work.

\textbf{Standard cells.} \textsc{ChipSAT} places only macros; standard cells are handled by the downstream analytical placer. This is consistent with the hierarchical placement paradigm used in industry, where macro placement is a separate upstream step.


%% ============================================================
\section{Conclusion}
\label{sec:conclusion}

We have identified a fundamental limitation of analytical macro placement: the features most predictive of post-route PPA are non-differentiable and therefore inaccessible to gradient-based optimizers. By introducing constraint programming (CP-SAT) to macro placement, we enable direct optimization of these features---minimum spacing, worst-case net span, density peaks, and routing congestion---with exact non-overlap enforcement.

Our framework, \textsc{ChipSAT}, embedded within ALNS, achieves the first simultaneous multi-metric PPA improvement on a ChiPBench circuit, reducing DRC by 26\% and TNS by 8\% on \texttt{swerv\_wrapper}. This result demonstrates that objective choice, not optimizer sophistication, is the primary bottleneck in macro placement quality.

The broader implication is that the field's decades-long focus on HPWL may be a consequence of tooling limitations rather than a principled design choice. Constraint programming offers a complementary paradigm that can optimize the objectives that actually matter for post-route PPA.


%%
%% Bibliography
\bibliographystyle{ACM-Reference-Format}
\bibliography{references}

\end{document}